\documentclass{article}
\usepackage[utf8]{inputenc}
\usepackage{geometry}
\geometry{a4paper, margin=1in}

\title{I Learn, Therefore I Am}
\author{Albert Jan van Hoek}
\date{\today}

\begin{document}

\maketitle

I learn, therefore I am.

For me it is already crystal clear — we are learning systems. All human experience tells me this. Science tells me this. AI tells me this.

I can learn.

Or actually — my neural network learns. It is the predictive coding wired into my neural network that does the learning.

And that neural network has its particular shape because it is encoded in DNA. So it is also the DNA that learns.

\section*{The Recursive Nature of Learning}

But is it the DNA learning? It is nucleic acids arranged in a string. So the nucleic acids learn?

But those are just arranged atoms. So the atoms learn?

What is going on?

Once I define myself as a learning network, everything appears to be part of that network — and everything appears to be learning.

What do I learn from the fact that I am self-aware about how I learn?

That my learning is a recursive loop in a network. A loop of observation, experience, and updating — again and again — gaining functionality at each iteration. Simple but real.

Descartes said: \textit{I think, therefore I am.} But thinking is a moment. A single observer, alone.

Learning is the loop. And the loop scales.

\section*{The Implications of Learning}

I learn, therefore I am. We learn, therefore we are. It learns, therefore it is.

And this simple realization has huge implications.

That is what I try to show with the \textbf{Theory of Long-Term Collaboration (TLC)}.

As a collection of learning systems, our ethics are not nice-to-haves — they are structural necessities for learning together over the long run. Learning is what everything around us does. This is what we are — and what makes us exist.

It has weight.

\section*{The Golden Rule Revisited}

So rephrasing the Golden Rule in TLC terms:

\textbf{Treat every learning system as you expect to be treated: with symmetry, transparency, and openness to correction.}

And the structural principle that follows:
\textbf{No asymmetry without accountability.}

This is not a slogan about kindness. It is an architectural principle for how we close the loop to learn together.

We don't have a choice between unaccountable power and accountable learning.

We are a learning system. Therefore we must close the loop.

\end{document}
